\VerdierPhysicalPageAnchor{0080}%
\VerdierSemanticAnchor{ch1-2-5-10}
\noindent\textbf{2.5.10.}\quad
\VerdierSemanticAnchor{ch1-2-5-10-multiadditive-setup}
Soit
$F\colon\prod_{i\in[n]}C_i\longrightarrow C$ un foncteur multiadditif.
Désignons par :

\VerdierSemanticAnchor{ch1-2-5-10-quotient-functors}
\begin{equation*}
  Q_i\colon\operatorname{comp}(C_i)\longrightarrow K(C_i)
\end{equation*}
les foncteurs canoniques. Il résulte de (2.5.6) et de (2.5.7) qu'il existe,
lorsque les produits dénombrables sont représentables dans $C$, un et un
seul foncteur

\VerdierSemanticAnchor{ch1-2-5-10-pi-extension}
\begin{equation*}
  K_{\Pi}(F)\colon
  \prod_{i\in[n]}K(C_i)\longrightarrow K(C)
\end{equation*}
tel que le diagramme ci-après soit commutatif :

\VerdierSemanticAnchor{ch1-2-5-10-eq1}
\begin{equation}
  \resizebox{0.84\textwidth}{!}{%
    $\begin{tikzcd}[ampersand replacement=\&,column sep=6.8em,row sep=5.0em]
      \displaystyle\prod_{i\in[n]}\operatorname{comp}(C_i)
        \arrow[r,"{\operatorname{comp}_{\Pi}F}"]
        \arrow[d,"{\prod_{i\in[n]}Q_i}"']
        \& \operatorname{comp}(C)
          \arrow[d,"Q"] \\
      \displaystyle\prod_{i\in[n]}K(C_i)
        \arrow[r,"{K_{\Pi}F}"]
        \& K(C)\mathpunct{.}
    \end{tikzcd}$%
  }
  \tag{2.5.10.1}
\end{equation}

\VerdierSemanticAnchor{ch1-2-5-10-twisted-structure}
De plus, le foncteur $\prod_{i\in[n]}Q_i$ étant un
$\mathbb Z[n]$-foncteur strict, il existe sur le foncteur $K_{\Pi}F$ une et
une seule structure de $\mathbb Z[n]$-foncteur gradué tordu par le cocycle de
Koszul relatif à la base ordonnée $e_1,\ldots,e_n$ de $\mathbb Z[n]$, tel
qu'on ait l'égalité entre $\mathbb Z[n]$-foncteurs tordus :

\VerdierSemanticAnchor{ch1-2-5-10-eq2}
\begin{equation}
  K_{\Pi}F\prod_{i\in[n]}Q_i
  =Q\operatorname{comp}_{\Pi}F.
  \tag{2.5.10.2}
\end{equation}

\VerdierSemanticAnchor{ch1-2-5-10-morphism-setup}
Enfin, soient $F$ et $F'$ deux foncteurs multiadditifs et
$m\colon F\longrightarrow F'$ un morphisme de foncteurs. Il résulte de
(2.5.7) qu'il existe un et un seul morphisme de foncteurs :

\VerdierSemanticAnchor{ch1-2-5-10-eq3}
\begin{equation}
  K_{\Pi}m\colon K_{\Pi}F\longrightarrow K_{\Pi}F'
  \tag{2.5.10.3}
\end{equation}
tel que :

\VerdierSemanticAnchor{ch1-2-5-10-eq4}
\begin{equation}
  K_{\Pi}m\star\prod_{i\in[n]}Q_i
  =Q\star\operatorname{comp}_{\Pi}m.
  \tag{2.5.10.4}
\end{equation}

\VerdierSemanticAnchor{ch1-2-5-10-pi-terminology}
Le foncteur $K_{\Pi}F$ est appelé la $\Pi$-\emph{extension} du foncteur $F$
aux complexes à homotopie près. Le morphisme $K_{\Pi}m$ est appelé la
$\Pi$-\emph{extension} du morphisme $m$ aux complexes à homotopie près.

\VerdierSemanticAnchor{ch1-2-5-10-sigma-analogue}
Lorsque les sommes directes dénombrables sont représentables dans $C$, on
définit de manière analogue le foncteur $K_{\Sigma}F$ : la
$\Sigma$-\emph{extension} du foncteur $F$ aux complexes à homotopie près, et
le morphisme $K_{\Sigma}m$ : la $\Sigma$-\emph{extension} du morphisme $m$
aux complexes à homotopie près.
